% !PRINTED 8
\weberpage{8}
Denkt man sich also $x_1, x_2, x_3, x_4$ in dieser Reihenfolge auf eine Kreisperipherie gesetzt, so erhält man immer dann ein positives echt gebrochenes $\varkappa^2$, wenn die $x_i$ so der Größe nach aufeinander folgen, daß man mit einem beliebigen als kleinstem anfängt und dann auf dem Kreise entweder nach rechts oder nach links weiter zählt. Wir drücken dies so aus:

\gesperrt{Damit $\varkappa^2$ ein positiver echter Bruch sei, müssen die Größen $x_1, x_2, x_3, x_4$ der Größe nach zyklisch aufeinander folgen.}

§ 3. \gesperrt{Lineare Transformation des elliptischen Differentials.}

Wenn sich die beiden Punkte $x_3$ und $x_4$ einem und demselben Punkte $x$ annähern, so nähern sich $x_3'$ und $x_4'$ dem entsprechenden Punkte $x'$ an. Wir setzen dann $(x_3x_4):(x_3'x_4') = dx:dx'$ und erhalten aus (3), § 2, die Beziehung zwischen den Differentialen $dx, dx'$:
\weq{1}{\frac{(x_1x_2)\,dx}{(x_1x)(x_2x)} = \frac{(x_1'x_2')\,dx'}{(x_1'x')(x_2'x')}.}

Ebenso ergibt sich, wenn man an Stelle von $x_1, x_2$ zwei andere Punkte $x_3, x_4$ setzt:
\weq{2}{\frac{(x_3x_4)\,dx}{(x_3x)(x_4x)} = \frac{(x_3'x_4')\,dx'}{(x_3'x')(x_4'x')},}

und wenn man multipliziert und die Wurzel zieht:
\weq{3}{\frac{\sqrt{(x_1x_2)(x_3x_4)}\,dx}{\sqrt{(x_1x)(x_2x)(x_3x)(x_4x)}} = \frac{\sqrt{(x_1'x_2')(x_3'x_4')}\,dx'}{\sqrt{(x_1'x')(x_2'x')(x_3'x')(x_4'x')}},}

\gesperrt{also eine Transformation des elliptischen Differentials durch eine lineare Substitution:}
\[
x = \frac{\alpha x' + \beta}{\gamma x' + \delta}.
\]

Diese Substitution ist vollkommen bestimmt, wenn zu drei beliebigen Punkten $x_1, x_2, x_3$ die zugehörigen Werte $x_1', x_2', x_3'$ willkürlich gegeben sind, und man erhält sie aus der Gleichheit des Doppelverhältnisses:
\weq{4}{\frac{(x_1x_2)(x_3x)}{(x_1x_3)(x_2x)} = \frac{(x_1'x_2')(x_3'x')}{(x_1'x_3')(x_2'x')},}

durch Auflösung nach $x$, und den vierten zu $x_4$ gehörigen Wert $x_4'$ erhält man aus
%CONT>>
